\documentclass[10pt]{article}
\usepackage[margin=20mm]{geometry}
\usepackage{amsmath,amssymb,mathtools,fontspec,unicode-math,xurl,hyperref}
\IfFontExistsTF{DejaVu Serif}{\setmainfont{DejaVu Serif}}{\setmainfont{Latin Modern Roman}}
\setmathfont{latinmodern-math.otf}
\setlength{\parindent}{0pt}\setlength{\parskip}{6pt}\sloppy
\title{A technical review of
Undecidability of translational monotilings}
\author{AI-assisted technical review}
\date{}
\begin{document}
\maketitle
Abstract: The exact manuscript arXiv:2309.09504v2 is reviewed. The reduction from domino constraints through decorated \(p\)-adic Sudoku to translational monotiling is traced, with attention to exceptional rows, precision, finite group encodings and periodic targets. A counterexample disproves the stated Tetris preservation corollary. Two further essential derivations use insufficient \(p\)-adic precision and omit cross-prime unit factors. These are assessed as Major at the level of the affected core derivations. Explicit restricted statements and corrected calculations recover the uses examined here; no counterexample to the main undecidability theorem is established. Additional local defects concern the claimed bijection, an unused topological interpretation and zero board width. The assessment applies to this manuscript, with scientific adjudication pending.\par
\section*{Exact source and review scope}
M25-0300 is Undecidability of translational monotilings, by Rachel Greenfeld and Terence Tao [paper]. The reviewed source is arXiv:2309.09504v2, submitted on 24 October 2023 at 19:41:53 UTC, with 44 pages and ten figures. Its 3,684,687-byte PDF has SHA256 `1c4aea6c36efb1f269f3df2e1bd7aec3a8b3bc6ad9cb72c1b7d0310bbad100ca`. The EMS record gives DOI 10.4171/JEMS/1673, submission on 24 October 2023, acceptance on 27 June 2024 and publication on 21 June 2025 [publisher]. The journal PDF was unavailable in the acquired material. No equivalence with that text is asserted. The frozen 2025 identifier and original title mapping are retained.\par
All 44 pages, including references, were read. Thirty-five source pages were actually viewed as images, including every figure and every page containing a substantive finding below. The small-prime illustrations are illustrative examples, not instances of the large-prime classification theorem. Exact integer calculations were executed independently, without running author code. The infinite Sudoku counterexamples are proved below for every required line; finite tests supplement those proofs and do not replace them.\par
The author publication entry, author discussion, exact arXiv history and EMS record were checked [author, discussion]. Bounded searches located no separate correction, supplement or associated code release. This is not an exhaustive absence assertion, and lack of code is not a mathematical defect. Two cited reduction papers were acquired at precisely identified versions [GT21, MSS22]. Their relevant statements and proof portions were read; neither entire reference was independently certified.\par
\section*{Quantifiers and dependency structure}
Theorem 1.3 concerns inputs consisting of a finite Abelian group \(G_0\), a subset \(E\subseteq G_0\), and a finite set \(F\subset\mathbb Z^2\times G_0\), and the existence of translations tiling \(\mathbb Z^2\times E\). The finite group varies with the input. Corollary 1.5 concerns a periodic target in \(\mathbb Z^d\), with \(d\) also part of the input. It does not establish undecidability for one fixed \(d\), or require the target to be the whole group. Remark 6.2 and Section 7 explicitly distinguish these stronger questions. The author discussion gives the same qualification.\par
The dependency chain is \[\begin{aligned}
 \text{domino constraints on }\mathbb N^2
 &\ \longrightarrow\ \text{decorated Sudoku}\\
 &\ \longrightarrow\ \text{weak functional constraints}\\
 &\ \longrightarrow\ \text{one translational prototile}.
\end{aligned}\] Section 2 uses the classical domino undecidability input and compactness. The finite-box argument in Lemma 2.5 is valid: compatible infinite configurations follow by a diagonal subsequence in a finite alphabet. This does not claim a computable bound on the box size witnessing failure. The conversion from Wang colors to neighbor relations is effective. The external undecidability theorem itself was not reproved.\par
Write \(B=\{0,\ldots,N-1\}\times\mathbb Z\). Definition 4.5 requires, for each nonvertical line \(m=jn+i\), a nondegenerate integer affine form \(an+b\) whose last nonzero base-\(p\) digit agrees with the line whenever \(\nu_p(an+b)\le1\). The function \(f_p\) has \(f_p(0)=1\). Its multiplicativity must not be used across a zero factor. Theorem 4.6 assumes \(p>48\), \(p^2\mid N\), and, in its converse, nonconstant columns. Its conclusion deliberately excludes the points where the limiting affine form is zero.\par
Theorem 4.7 supplies an initial affine description outside one exceptional residue set. Theorem 4.9 iterates that description to arbitrary finite precision, using the Tetris map \(m\mapsto pm\). Compactness then gives Theorem 4.6(ii). In Section 5 the fixed choices are \[p_1=53,\quad p_2=59,\quad N=p_1^2p_2^2=9,778,129.\] The symbols decorating the two digit sequences encode neighboring domino values at adjacent valuation levels. Thus the precision and coefficient comparisons in that section are essential to the reduction, even though the coefficients themselves are auxiliary objects.\par
\section*{Confirmed findings}
Problem 1 (Major: Corollary 4.10 does not preserve nonconstant columns).  On p. 25, Corollary 4.10 asserts that every \(S_{p,N}\)-Sudoku solution with nonconstant columns retains that property after \(F(n,m)\mapsto F(n,pm)\). Take \(p=53\), \(N=p^2\), and \[F(n,m)=f_p(m+1).\] All stated prime and width hypotheses hold. Every column is nonconstant, since its values at \(m=0,1\) are \(1,2\).\par
Here is a direct verification of every line constraint. For a line with \(j=i+1=0\), take the constant affine form \(1\). Otherwise put \[d=\min\{\nu_p(j),\nu_p(i+1)\},\qquad
 a=j/p^d,\quad b=(i+1)/p^d.\] The pair \((a,b)\) is nondegenerate modulo \(p\). If \(an+b\ne0\), removal of the power \(p^d\) preserves its last nonzero digit. If \(an+b=0\), both original and normalized arguments are zero and both have the assigned value \(1\). Consequently \(F(n,jn+i)=f_p(an+b)\) everywhere on the line, in particular at every point required by Definition 4.5. Thus this is a genuine admissible Sudoku solution, not merely a sampled array.\par
Its Tetris image is \[F(n,pm)=f_p(pm+1)=1
                         \qquad(m\in\mathbb Z).\] Every column becomes constant. This disproves the preservation claim under all its printed hypotheses. The independent calculation checked 151,650 regulated line digits on selected lines and 6,001 Tetris values; the proof above establishes the unrestricted statement.\par
The error is in the corollary’s inference from the \(r=1\) affine description. A coefficient \(B^{(1)}\) that is a unit does not suffice: the constant \(A^{(1)}n+C^{(1)}\) can keep the entire sampled column away from the singular residue class. In this example it is \(1\). Proposition 3.4(iii) preserves being a Sudoku solution; it does not supply the additional column condition.\par
A sufficient corrected hypothesis is the normalization already used in (4.11): \[F(n,m)=\pi_p(bm)\quad\text{when }p\nmid m,\qquad p\nmid b.\] For such \(F\), apply Theorem 4.7 to its Tetris image, obtaining coefficients \(A_1,B_1,C_1\). The comparison of affine forms on pp. 23–24, culminating in (4.22), gives \(B_1\equiv b\pmod p\). This comparison precedes and does not invoke Corollary 4.10. Since \(B_1\) is a unit, Theorem 4.7 implies that the Tetris image has nonconstant columns. Hence the restricted claim is proved.\par
The general induction on p. 25 first imposes (4.23), which includes this normalization, before invoking the corollary. With the precision repair in Problem 3, that particular invocation is justified. The false unrestricted corollary remains false and must be replaced. Major severity concerns a named preservation guarantee in the core structural recursion; the fact that its actual later use admits a verified restriction is recorded separately, and is not a counterexample to Theorem 4.6 or Theorem 1.3.\par
Problem 2 (Minor: the claimed bijection loses the exceptional row).  Remark 4.11, pp. 26–27, upgrades the description \[F(n,m)=c f_p(m-Dn-E)\quad(m\ne Dn+E)\] to a one-to-one correspondence with triples \((c,D,E)\). For the same \(p=53\) and \(N=p^2\), define, for \(k=1,2\), \[F_k(n,m)=
 \begin{cases}f_p(m),&m\ne0,\\k,&m=0.\end{cases}\] Both have nonconstant columns. On the zero row use the constant nondegenerate form \(k\). On every other line, divide \(jn+i\) by the common power of \(p\) in \(j,i\). Any zero of the resulting nondegenerate form has infinite valuation and is excluded from the constraints in Definition 4.5. All regulated points therefore agree. Both \(F_k\) are admissible, but they are distinct and both satisfy the displayed description with \((c,D,E)=(1,0,0)\). This refutes the asserted bijection. It does not refute the root-excluding conclusion of Theorem 4.6(ii).\par
Two printed identities in the accompanying uniqueness argument also fail. With \(F=f_p(m)\), \(m=p-1=52\), both \(m\) and \(m+1\) are nonzero, but in \(\mathbb F_p\) \[F(n,53)-F(n,52)=1-52=2\ne1=c.\] The increment identity is valid if both affine arguments are units, not merely nonzero. In the next display on p. 27, the definition \(h_{n,m}=(m-Dn-E)/p^{r-1}\) gives \(c\pi_p(h+h_{n,m})\), valid when \(h\not\equiv-h_{n,m}\pmod p\), rather than the printed difference. For \(r=1,D=E=0,m=1,h=2,c=1\), the actual value is \(3\); the printed value is \(1\), although its exclusion condition is met.\par
Uniqueness of the parameters can be retained independently. Comparing two descriptions away from their exceptional residues gives equality of their nonzero affine forms modulo \(p\), which first determines \(c\) and \(Dn+E\bmod p\). Using columns \(n=0,1\) determines \(D,E\bmod p\). At each successive precision, restrict \(m\) to the already determined root residue and compare the two affine functions of the next digit, away from at most two excluded residues. This determines the next digits of \(D,E\). This argument proves uniqueness of the parameters from \(F\); it does not determine \(F\) on an actual zero locus from the parameters.\par
The last display on p. 17 has a related zero-domain issue: with \(A=C=0,B=2,m=0\), its left side is \(f_p(0)=1\) and its right side is \(f_p(2)f_p(0)=2\). Restricting that factorization to nonzero arguments is sufficient for the nonconstant-column conclusion, since the unit residues already give varying values. Section 5 uses existence of the root-excluding representation, not the false bijection. These local extra assertions therefore do not remove a needed core input.\par
Problem 3 (Major: the precision used for last-digit normalization is insufficient).  On p. 25, immediately before (4.23), the affine forms are matched modulo \(p^r\), after which equality of \(f_p\) is asserted through valuation \(r\). On p. 30, the CRT choices of \(D,E\) likewise use \(p_l^{r_l}\), followed by (5.1) through valuation \(r_l\). The required modulus is one power higher.\par
An exact test of the inference is \(F(n,m)=f_p(m+p^r)\), with affine coefficients \(A=0,B=1,C=p^r\), and \(D=E=0\). The printed congruence modulo \(p^r\) is satisfied. At \(m=p^r\), whose valuation is \(r\), the original value is \(2\), whereas the claimed normalized value is \(1\). This is an admissible nonconstant-column Sudoku solution by the line argument in Problem 1 with a different shift. The independent script verifies this discrepancy for \(r=1,2,3,4\); the displayed calculation proves it for every \(r\ge1\).\par
This test disproves the asserted consequence of the stated congruence. It does not disprove the existence of better choices of \(D,E\), or the existence conclusions of the theorems. For a unit \(B\), choose \[D\equiv-A/B,\qquad E\equiv-C/B\pmod {p^{r+1}}.\] After shearing, the two affine arguments differ by a multiple of \(p^{r+1}\). Whenever one has valuation \(k\le r\), division by \(p^k\) shows that their unit residues modulo \(p\), and therefore their \(f_p\) values, agree. For the two-prime argument choose \(D,E\) simultaneously modulo \(p_1^{r_1+1}\) and \(p_2^{r_2+1}\). CRT provides these choices because the primes are distinct. No enlargement of the board or new hypothesis is needed.\par
This is an invalid precision inference at every recursive scale and at the normalization feeding domino recovery. Its placement in the core derivations warrants Major severity at that level, even though the stronger congruence supplies an explicit repair. The repair, rather than the insufficient printed modulus, is used in the subsequent checks in this report.\par
Problem 4 (Major: the domino-recovery calculation omits cross-prime units).  At the bottom of p. 31, after setting \(j=p_1^{s_1}p_2^{s_2}\), the proof asserts \(F_l(n,jn)=\pi_{p_l}(cn)\) for \(p_l\nmid n\). The correct coefficient is \[\lambda_1=\pi_{p_1}(c p_2^{s_2}),\qquad
 \lambda_2=\pi_{p_2}(c p_1^{s_1}).\] Removing the power of one prime leaves the other prime as a unit.\par
For a complete admissible instance, take \(W=\{w\}\), both domino relations equal to \(\{(w,w)\}\), and the constant decoration \(w\). The forward construction of Proposition 5.2 gives \[F(n,m)=(f_{53}(m),f_{59}(m),w).\] All line witnesses follow by division at each prime and CRT modulo \(53^2,59^2\), as in the forward proof. Both digit components have nonconstant columns. Use \(r=(1,1)\), \(s=(0,1)\), \(j=59\), \(c=1\), and \(n=1\). These lie in the prescribed ranges, but \[F_1(1,59)=6\ne1.\] The same missing factor occurs in the level-one comparison: \(F_1(53,59\cdot53)=6\), rather than \(1\). The exact script checks both valuation levels for several \(s,n\). Thus the coefficient identity is false for a genuine decorated Sudoku solution satisfying the operative hypotheses.\par
Here is the repaired domino extraction. For each \(l\), compare \(\lambda_l n\) with the line witness \(a_{j,0}n+b_{j,0}\) away from the excluded residue classes. Affine extrapolation gives \[a_{j,0}\equiv\lambda_l\pmod {p_l},\qquad
 b_{j,0}\equiv0\pmod {p_l}.\] Both \(\lambda_l\) are units. Consequently (5.6) still holds when \(n\) is a unit at both primes. Multiplication by these unit factors permutes the pairs of nonzero digit residues, so CRT still eliminates the final-digit dependence in (5.5).\par
For a horizontal domino with \(s_1+1\le r_1\), compare the first component at the column \(p_1n\), using \(0\le n<N/p_1\) and \(p_1\nmid n\). The two expressions are \[\lambda_1 n,\qquad
 \pi_{p_1}\!\left(a_{j,0}n+b_{j,0}/p_1\right).\] Extrapolation now yields \(b_{j,0}\equiv0\pmod {p_1^2}\). At an \(n\) coprime to both primes, the witness therefore has valuation pair \((1,0)\) at \(p_1n\) and \((0,0)\) at \(n\). Equation (5.3) reads the two entries of \(T_{j,0}\), giving the required horizontal relation. The vertical argument is symmetric. There are more than enough residue representatives in the stated width. Using \(0\le n<N/p_1\) also removes the printed upper endpoint that would put column \(p_1n=N\) outside the board.\par
The preceding sentence on p. 31 asserting an equivalence between \(\nu_{p_l}(j)>\nu_{p_l}(i)\) and equality of last digits also needs the equality of the two finite valuations. That equality is explicit in the actual application immediately below it, so the application remains valid.\par
The missing factors invalidate the printed coefficient comparison on which the reverse domino reduction depends. Major severity concerns that core calculation. The corrected unit factors recover the valuation transitions checked above; the example does not disprove the equivalence of domino and decorated-Sudoku solvability.\par
Problem 5 (Minor: the sets in Remark 2.4 do not form a topology).  On p. 8, the proposed open sets \(U_\Omega\) contain the vertices of \(\Omega\) and all domino edges whose endpoints lie in \(\Omega\). For adjacent vertices \(s,t\), the sets \(U_{\{s\}}=\{s\}\) and \(U_{\{t\}}=\{t\}\) belong to the specified family. Their union \(\{s,t\}\) does not: its vertex set forces \(\Omega=\{s,t\}\), for which \(U_\Omega\) also contains the edge \([s,t]\). Closure under unions fails.\par
A corrected incidence topology permits a vertex set \(S\) together with any edge set \(D\) whose endpoints belong to \(S\). Sections assign labels to \(S\) and impose the relations for the edges in \(D\). Compatible sections glue, because each included edge and its endpoints occur together in a member of an open cover. On the original \(U_\Omega\) this gives the intended domino sections. Vertex stalks have alphabet \(W\); edge stalks instead have the allowed pairs in the appropriate relation, so the assertion on p. 9 that all stalks are \(W\) also needs correction. The manuscript explicitly does not use this sheaf interpretation later. This is a substantive local mathematical error, with no core reduction depending on it.\par
Problem 6 (Minor: zero board width contradicts the uniqueness assertion).  Section 1.5, p. 6, explicitly includes zero in \(\mathbb N\). Definition 4.5 allows a natural width, and Theorem 4.7, p. 18, only requires \(p^2\mid N\). At \(N=0\), the board is empty. The unique empty function satisfies all line conditions, and the affine identity (4.5) is vacuous. Both constant forms \(1\) and \(2\) are nondegenerate modulo \(p=53\) and satisfy it, but their coefficients are different modulo \(p\), contradicting the stated uniqueness. Lemma 4.8 also cannot locate a square on that board. State \(N>0\), equivalently \(N\ge p^2\) here. Every main application already uses the positive width \(53^2 59^2\); the empty-board defect has no effect on that reduction.\par
\section*{Further checks and interpretive boundaries}
\subsection*{The structural recursion after the stated repairs}
The good-square counting argument uses the coefficient \(48/p<1\), with the remaining finite-width term harmless for \(N\ge p^2\). The row, diagonal and antidiagonal propagation on a \(4\)-by-\(4\) square was checked against the figure and its affine identities. Affine extrapolation removes at most two exceptional residue classes, leaving distinct usable residues since \(p>48\). In the \(r=1\) comparison, the relevant arithmetic progression has step coprime to \(p\) and contains at least \(p\) terms. It therefore provides every residue needed to derive \(B_1\equiv B^{(0)}\pmod p\).\par
For the general induction use the stronger precision in Problem 3 and the normalized Tetris claim in Problem 1. At the level \(p^r m\), factor the new affine form by its common power \(p^j\), with \(j\le r-1\). Comparing its unit residues with \(B^{(r)}m\) outside at most two exceptional affine zero sets forces \(j=r-1\), \(A_{r+1},C_{r+1}\) divisible by \(p^r\), and \(B_{r+1}\equiv B^{(r)}\pmod p\). The coefficients \((pA_{r+1},B_{r+1},pC_{r+1})\) then cover the next valuation level. The first display on p. 25 should use \(m'\), rather than \(m\), inside \(A_1n+B_1m'+C_1\) when \(m=pm'\).\par
Compactness in \(\mathbb Z_p^3\) gives a convergent subsequence, with the middle coefficient still a unit because the units form a closed set. At any nonzero limiting affine argument, its valuation and last digit eventually stabilize. This proves precisely the root-excluding limit assertion. It imposes no value at an actual zero and therefore cannot justify the bijection refuted above.\par
\subsection*{Effective finite rules and the encoding}
The decorated rule is finite and effectively enumerable despite the large board. For a line witness, its relevant valuation and digit tests depend only on the coefficients modulo \(p_1^2p_2^2\); the possible small domino decorations are finite. The forward construction divides the two line coefficients by their common \(p_l\)-power separately, uses CRT, and translates the domino configuration by the two removed valuations. The identically zero line is handled separately. No running-time bound or explicit enumeration of all \(9,778,129\) columns was claimed or attempted in this review.\par
The choice \(q=p_1p_2\) in the initial condition requires genuine permutations of \(\mathbb Z/q\mathbb Z\). On the converse side, a CRT linear permutation with unit slope supplies the needed labels. In Section 6, finite conjunctions and existentially quantified auxiliary functions have finite constructive descriptions. The Boolean constraint uses an order-two direction; the non-wraparound bound on \(L\) is essential. The auxiliary-sum construction of Lemma 6.12 works for odd numbers of coordinates by excluding the two constant sign vectors; an extension can be chosen on each order-two orbit and negated on its partner. Exact checks for \(3,5,7,9\) coordinates covered 680 sign vectors. These are checks of the elementary arithmetic, not exhaustive tests of all encodings. In Lemma 6.13, consecutive windows each contain every antipodal label once. Comparing adjacent windows makes the antipodal label \(q\)-periodic and hence a permutation. The signs themselves need not be periodic.\par
There are several typing and indexing omissions that should be corrected before treating Section 6 as a literal specification. Use columns \(0,\ldots,N-1\) consistently, scalar sign coordinates \(\omega_{a,b,n}\in\{-1,1\}\), and codomain \(\Sigma\) for \(F\). Proposition 6.3 should say that \(F\) tiles \(G\times E'\); its proof has that target. Most significantly, add the word permutation to clause (4) on p. 41. The proof on p. 42 and Lemma 6.13 explicitly use that intended meaning.\par
For clarity, dropping bijectivity really would change the constraint. After consistent reindexing, take \(N=1,q=2\), \(\Sigma=\{a\}\), the single constant Sudoku word, and \(C=\{(0,a)\}\). No permutation can satisfy the initial condition. With \(s_0=2,L=13\), \[\iota_0(0)=(1,1),\quad \iota_0(1)=(1,-1),\quad
 \iota_1(a)=(1,1),\] all four scalar functions \(\alpha_{a,b,0}((i,j),t)=(-1)^t\) satisfy (1)–(3) and the weakened (4) with the constant map \(\sigma\equiv0\). The dimension and \(L>4s_0N+4\) conditions hold. This is a test of the weakened reading after reindexing; it is not a counterexample to the intended bijective Theorem 6.4, whose surrounding text states the missing requirement. It is recorded as a clarification, not an additional Major finding.\par
With that meaning explicit, invariance along \((-n,1)\) makes the decoded functions depend only on \(jn+i\). The disjoint antipodal images make the decoded symbols unique, and the permutation condition gives the original initial condition. The simultaneous global sign in \(\Omega\cup-\Omega\) is necessary; independent signs would not encode the same word.\par
\section*{External reduction checks}
For the merger into one tiling equation, Theorem 1.15 of the exact GT21 version [GT21] applies with one translation unknown. Its hypotheses include nonempty tiles and an extra cyclic group of order greater than the number of equations. The proof on pp. 15–16 uses the fact that a proper consecutive interval in that cyclic group has trivial translation stabilizer. Thus all centers have the distinguished extra coordinate, and slicing recovers the original equations. This keeps one prototile. The independent script checked that stabilizer fact in 300 small cases; the interval argument proves the general fact.\par
The functional graph equation in the manuscript has the correct shift sign: a tile at \(-h\) meets a fiber at \(x\) through \(\alpha(x+h)\). Empty or contradictory constraints must be preprocessed when invoking the nonempty-tile merger. An empty tile with empty target is discarded; an empty tile with nonempty target, or a nonempty tile with empty target under the graph condition, is an immediate impossibility. One can output a fixed impossible periodic-target instance in those cases. Overlapping components at the same shift likewise signal a violated disjoint-sum condition and can be detected before the merger.\par
Theorem 1.1 of MSS22 [MSS22] is stated for tiling the entire quotient group. Its printed statement alone should not be treated as a proved periodic-subset variant. For the actual constraints here, the GT21 periodic-target construction provides a route. The following explanation is an applicability check using its rigid-tile lemma, rather than a claim to have reproved that lemma.\par
Lift the function domain \(\mathbb Z^2\times\mathbb Z/2\mathbb Z\) to \(\mathbb Z^3\) and impose period two in its last coordinate. The encoded codomain is \(H=(\mathbb Z/L\mathbb Z)^K\), with \(L\ge5\). Let \(\pi:\mathbb Z^K\to H\), \(\Lambda=\ker\pi\), and take the finite rigid fundamental tile \(R\) from GT21 Lemma 9.3. For a center set \(A\subset\mathbb Z^3\times\mathbb Z^K\), the equation \[A\oplus(\{0\}\times R)=\mathbb Z^{3+K}\] forces each fiber \(A_x\) to be a coset \(\pi^{-1}(\alpha(x))\). For each finite functional rule, replace a subset \(E_{ij}\subset H\) by its representatives \(\widetilde E_{ij}=R\cap\pi^{-1}(E_{ij})\), and form the finite tile \[F_i=\bigsqcup_j \{-h_{ij}\}\times\widetilde E_{ij}.\] Then \[A\oplus F_i=\mathbb Z^3\times\pi^{-1}(E_i')\] is exactly the corresponding disjoint functional partition. The target is periodic. Theorem 9.1 of GT21 combines this finite system, still with one center unknown, into a single tile in one additional integer dimension and a periodic target. The dimension varies with the finite encoding, as required.\par
GT21 Lemma 9.3 and its geometric argument were inspected on pp. 27–28, including the bump-and-hole figure. Its asserted exhaustive geometric case analysis was not independently enumerated in every dimension; this remains an external proof dependency. The full MSS22 proof, the source’s suggested connected Euclidean analogue, and all later variants were not independently certified. Logical independence statements also retain the set-theoretic assumptions identified in GT21, rather than asserting consistency of ZFC from within ZFC.\par
\section*{Assessment}
Problems 1, 3 and 4 identify failed core guarantees or calculations, with exact counterexamples to the particular assertions and explicit dependencies. They warrant an overall Major assessment under a claim-level standard. The repaired normalization, restricted Tetris step and cross-prime calculations recover the portions of the main reduction checked here. No contradiction to the main undecidability conclusion is established. Problems 2, 5 and 6 are local defects with their narrower impact stated.\par
This distinction is material: a repairable core derivation is still a failed derivation as printed, while a counterexample to an auxiliary assertion must not be represented as a counterexample to the entire paper. The report does not certify the unavailable journal version or every external theorem. It is an initial technical review; scientific decision remains pending, human reviews are empty, and no independent acceptance is claimed.\par
9 R. Greenfeld and T. Tao, Undecidability of translational monotilings, arXiv:2309.09504v2, 24 October 2023, https://arxiv.org/abs/2309.09504v2. EMS Press, article record, DOI 10.4171/JEMS/1673, https://ems.press/journals/jems/articles/14298908 (accessed 27 October 2026). T. Tao, publication list, https://teorth.github.io/tao-web/papers.html (accessed 27 October 2026). T. Tao, author discussion of the paper, 18 September 2023, https://terrytao.wordpress.com/2023/09/18/undecidability-of-translational-monotilings/. R. Greenfeld and T. Tao, Undecidable translational tilings with only two tiles, or one nonabelian tile, arXiv:2108.07902v2, 23 August 2022, 48 pp., https://arxiv.org/abs/2108.07902v2. T. Meyerovitch, S. Sanadhya and Y. Solomon, A note on reduction of tiling problems, arXiv:2211.07140v1, 14 November 2022, 10 pp., https://arxiv.org/abs/2211.07140v1.\par
\end{document}
